\textbf{Design.} An open-weight code model writes MySQL for 40 held-out analyst questions. The gold SQL comes from
the OLAP pack and variants of it (\tt{text2sql/questions\_test.py}). Each question states the columns it expects,
so ``same rows'' is well defined. Matching is order- and column-order-insensitive, with numbers compared at two
decimals. 22 questions can be answered exactly from an aggregate table, and they define the
\emph{aggregate-navigation rate}. The three conditions are:
\begin{itemize}
\item \textbf{No schema.} The model gets only the dialect and the question.
\item \textbf{Schema in prompt.} The context is built \emph{programmatically from the live warehouse}
(script \tt{schema\_context.py}): compact DDL with keys and foreign keys for the 10 dimension, fact and
aggregate tables, each table's grain statement read back from its comment, and a data dictionary of every
low-cardinality value. With the instructions, the prompt is 2,438 tokens. Staging tables are deliberately excluded.
\item \textbf{LoRA + schema.} The same prompt, after LoRA fine-tuning [5] on 200 (question, SQL) pairs. The pairs
come from 21 parameterized templates (\tt{text2sql/make\_train\_pairs.py}). Each pair was executed and checked to
return rows, and none duplicates a test question or its gold SQL.
\end{itemize}

\textbf{What ran where.} The proposal's degradation path was used for the prompt-only conditions.
Qwen2.5-Coder-7B-Instruct [6], the proposal's first-choice model, ran as a 4-bit GGUF (Q4\_K\_M) through
\tt{llama.cpp} on the 2-vCPU build machine, with greedy decoding. The 1.5B sibling ran the same way as a size
contrast. The LoRA condition needs a GPU. It is packaged for the course's NRP hub (\tt{text2sql/nrp/}): training
script, the identical prompt, a generate-only evaluator, and a scorer that runs back on the warehouse machine.
Its rows in Table~\ref{tab:llm} are filled from that run. No LoRA number is reported until it exists.

\begin{table}[H]\centering\small
\caption{Text-to-SQL results on the 40 held-out questions (22 aggregate-eligible). Execution match = same multiset of rows as the gold query; aggregate navigation = share of aggregate-eligible questions answered from an \tt{agg\_} table.}\label{tab:llm}
\begin{tabular}{@{}lllrrrr@{}}
\toprule
\textbf{Model} & \textbf{Condition} & \textbf{Run} & \textbf{Exec.\ match} & \textbf{Invalid SQL} & \textbf{Agg.\ navigation} & \textbf{s/q} \\ \midrule
Qwen2.5-Coder-7B (4-bit) & no schema & CPU & 0.0\% (0/40) & 100.0\% & 0.0\% & 29 \\
Qwen2.5-Coder-7B (4-bit) & schema in prompt & CPU & 47.5\% (19/40) & 22.5\% & 90.9\% & 57 \\
\addlinespace[2pt]
Qwen2.5-Coder-1.5B (4-bit) & no schema & CPU & 0.0\% (0/40) & 100.0\% & 0.0\% & 6 \\
Qwen2.5-Coder-1.5B (4-bit) & schema in prompt & CPU & 15.0\% (6/40) & 75.0\% & 81.8\% & 12 \\
\addlinespace[2pt]
Qwen2.5-Coder-1.5B (fp16) & no schema & NRP GPU & \multicolumn{4}{c}{\pending} \\
Qwen2.5-Coder-1.5B (fp16) & schema in prompt & NRP GPU & \multicolumn{4}{c}{\pending} \\
Qwen2.5-Coder-1.5B (fp16) & LoRA + schema & NRP GPU & \multicolumn{4}{c}{\pending} \\
\bottomrule
\end{tabular}
\end{table}


\textbf{Results.} Schema context is decisive. Without it, neither model produced a single executable query. Every
attempt invented plausible tables (\tt{oulad\_dw.vle}, \tt{students}, \tt{enrollments}). With the
warehouse-generated context, the 7B model matched \textbf{19 of 40} questions (47.5\%), and only 9 of its queries
failed to execute. The 1.5B model matched 6 of 40, and 75\% of its SQL was invalid. On aggregate navigation, the
7B model reached for an \tt{agg\_} table on 20 of the 22 eligible questions (\textbf{91\%}), scanned the base fact
on only 2, and every one of its 8 correct aggregate-eligible answers used an aggregate. The grain comments make
the aggregates \emph{discoverable}: the model knows they exist and when to prefer them. Discoverable is not the
same as \emph{usable}, though. The 7B model was more accurate on the questions that did \emph{not} involve
aggregates (61\%, 11/18) than on those that did (36\%, 8/22).

\textbf{Where the design was not legible.} Each of the 7B model's 21 failures was classified by primary cause.
\begin{itemize}
\item \textbf{Attribute on the wrong table (9; all of the invalid SQL).} For example, \tt{code\_presentation} on
an aggregate, \tt{highest\_education} on the enrollment, or \tt{pres\_year} on the enrollment. The model assumes
denormalized attributes that a strict star schema keeps in the dimensions.
\item \textbf{Surrogate key used as a value (4).} For example, \tt{week\_key = 2} and \tt{week\_key <= 0}. The
surrogate key is offset from \tt{week\_num} by 4, so these queries run but return the wrong week.
\item \textbf{Composite code in the wrong column (2 primary, 7 overall).} For example,
\tt{code\_presentation = 'CCC-2014J'} where \tt{module\_presentation} was meant.
\item \textbf{Wrong grain in the arithmetic (6).} Averaging over enrollment-\emph{weeks} instead of students, and
in one case joining the presentation-level aggregate to enrollments, which multiplies the clicks.
\end{itemize}
The second and fourth causes are specific to dimensional design. Integer surrogate keys and pre-aggregated grains
are exactly what make the warehouse fast and correct for its designer, and exactly what an outside reader
misreads. The result suggests three inexpensive design changes:
\begin{itemize}
\item Carry \tt{week\_num} and \tt{module\_presentation} as degenerate attributes on the aggregates.
\item State in the dictionary that \tt{week\_key} is not a week number.
\item Include one worked aggregate query in the context.
\end{itemize}
Testing those changes is future work. The LoRA condition tests the other route, which is to teach the model the
schema's conventions rather than to change the schema.
